%% kapa: exam=1 round=20 year=2009 period=2 form=A subject=appraisal_law subject_ko=부동산관계법규 seq=33 number=73
%% source: Q-Net 감정평가사 기출문제 (cst00309) artlSeq=5024178 file=extracted/1008/5024178/2009년 제20회 감정평가사 1차시험 기출문제/2교시.hwp
%% conversion: kordoc 4.24.51 -> tools/gen1.py (auto); review=proofread-sonnet; flags=none
\begin{question}{73}
지적법령상 면적에 관한 설명으로 옳지 \uline{않은} 것은?
\begin{choices}
  \item 면적을 측정하는 경우 도곽선의 길이에 0.5밀리미터 이상의 신축이 있는 때에는 이를 보정하여야 한다.
  \item 면적이라 함은 지적공부에 등록한 필지의 수평면상 넓이를 말하며, 단위는 제곱미터로 한다.
  \item 토지의 면적결정을 함에 있어서 면적에 제곱미터 미만의 끝수가 있는 경우 0.5제곱미터를 초과하는 때에는 올리는 것이 원칙이다.
  \item 축척 600분의 1 지적도 시행지역에서 측정한 1필지의 면적이 0.09제곱미터인 경우에는 1제곱미터로 등록한다.
  \item 경위의측량방법에 의해 세부측량을 시행한 지역의 필지별 면적측정방법은 경계점좌표에 의한다.
\end{choices}
\end{question}
\answer{4}
